% Options for packages loaded elsewhere
% Options for packages loaded elsewhere
\PassOptionsToPackage{unicode}{hyperref}
\PassOptionsToPackage{hyphens}{url}
\PassOptionsToPackage{dvipsnames,svgnames,x11names}{xcolor}
%
\documentclass[
  letterpaper,
  DIV=11,
  numbers=noendperiod]{scrartcl}
\usepackage{xcolor}
\usepackage{amsmath,amssymb}
\setcounter{secnumdepth}{-\maxdimen} % remove section numbering
\usepackage{iftex}
\ifPDFTeX
  \usepackage[T1]{fontenc}
  \usepackage[utf8]{inputenc}
  \usepackage{textcomp} % provide euro and other symbols
\else % if luatex or xetex
  \usepackage{unicode-math} % this also loads fontspec
  \defaultfontfeatures{Scale=MatchLowercase}
  \defaultfontfeatures[\rmfamily]{Ligatures=TeX,Scale=1}
\fi
\usepackage{lmodern}
\ifPDFTeX\else
  % xetex/luatex font selection
\fi
% Use upquote if available, for straight quotes in verbatim environments
\IfFileExists{upquote.sty}{\usepackage{upquote}}{}
\IfFileExists{microtype.sty}{% use microtype if available
  \usepackage[]{microtype}
  \UseMicrotypeSet[protrusion]{basicmath} % disable protrusion for tt fonts
}{}
\makeatletter
\@ifundefined{KOMAClassName}{% if non-KOMA class
  \IfFileExists{parskip.sty}{%
    \usepackage{parskip}
  }{% else
    \setlength{\parindent}{0pt}
    \setlength{\parskip}{6pt plus 2pt minus 1pt}}
}{% if KOMA class
  \KOMAoptions{parskip=half}}
\makeatother
% Make \paragraph and \subparagraph free-standing
\makeatletter
\ifx\paragraph\undefined\else
  \let\oldparagraph\paragraph
  \renewcommand{\paragraph}{
    \@ifstar
      \xxxParagraphStar
      \xxxParagraphNoStar
  }
  \newcommand{\xxxParagraphStar}[1]{\oldparagraph*{#1}\mbox{}}
  \newcommand{\xxxParagraphNoStar}[1]{\oldparagraph{#1}\mbox{}}
\fi
\ifx\subparagraph\undefined\else
  \let\oldsubparagraph\subparagraph
  \renewcommand{\subparagraph}{
    \@ifstar
      \xxxSubParagraphStar
      \xxxSubParagraphNoStar
  }
  \newcommand{\xxxSubParagraphStar}[1]{\oldsubparagraph*{#1}\mbox{}}
  \newcommand{\xxxSubParagraphNoStar}[1]{\oldsubparagraph{#1}\mbox{}}
\fi
\makeatother

\usepackage{color}
\usepackage{fancyvrb}
\newcommand{\VerbBar}{|}
\newcommand{\VERB}{\Verb[commandchars=\\\{\}]}
\DefineVerbatimEnvironment{Highlighting}{Verbatim}{commandchars=\\\{\}}
% Add ',fontsize=\small' for more characters per line
\usepackage{framed}
\definecolor{shadecolor}{RGB}{241,243,245}
\newenvironment{Shaded}{\begin{snugshade}}{\end{snugshade}}
\newcommand{\AlertTok}[1]{\textcolor[rgb]{0.68,0.00,0.00}{#1}}
\newcommand{\AnnotationTok}[1]{\textcolor[rgb]{0.37,0.37,0.37}{#1}}
\newcommand{\AttributeTok}[1]{\textcolor[rgb]{0.40,0.45,0.13}{#1}}
\newcommand{\BaseNTok}[1]{\textcolor[rgb]{0.68,0.00,0.00}{#1}}
\newcommand{\BuiltInTok}[1]{\textcolor[rgb]{0.00,0.23,0.31}{#1}}
\newcommand{\CharTok}[1]{\textcolor[rgb]{0.13,0.47,0.30}{#1}}
\newcommand{\CommentTok}[1]{\textcolor[rgb]{0.37,0.37,0.37}{#1}}
\newcommand{\CommentVarTok}[1]{\textcolor[rgb]{0.37,0.37,0.37}{\textit{#1}}}
\newcommand{\ConstantTok}[1]{\textcolor[rgb]{0.56,0.35,0.01}{#1}}
\newcommand{\ControlFlowTok}[1]{\textcolor[rgb]{0.00,0.23,0.31}{\textbf{#1}}}
\newcommand{\DataTypeTok}[1]{\textcolor[rgb]{0.68,0.00,0.00}{#1}}
\newcommand{\DecValTok}[1]{\textcolor[rgb]{0.68,0.00,0.00}{#1}}
\newcommand{\DocumentationTok}[1]{\textcolor[rgb]{0.37,0.37,0.37}{\textit{#1}}}
\newcommand{\ErrorTok}[1]{\textcolor[rgb]{0.68,0.00,0.00}{#1}}
\newcommand{\ExtensionTok}[1]{\textcolor[rgb]{0.00,0.23,0.31}{#1}}
\newcommand{\FloatTok}[1]{\textcolor[rgb]{0.68,0.00,0.00}{#1}}
\newcommand{\FunctionTok}[1]{\textcolor[rgb]{0.28,0.35,0.67}{#1}}
\newcommand{\ImportTok}[1]{\textcolor[rgb]{0.00,0.46,0.62}{#1}}
\newcommand{\InformationTok}[1]{\textcolor[rgb]{0.37,0.37,0.37}{#1}}
\newcommand{\KeywordTok}[1]{\textcolor[rgb]{0.00,0.23,0.31}{\textbf{#1}}}
\newcommand{\NormalTok}[1]{\textcolor[rgb]{0.00,0.23,0.31}{#1}}
\newcommand{\OperatorTok}[1]{\textcolor[rgb]{0.37,0.37,0.37}{#1}}
\newcommand{\OtherTok}[1]{\textcolor[rgb]{0.00,0.23,0.31}{#1}}
\newcommand{\PreprocessorTok}[1]{\textcolor[rgb]{0.68,0.00,0.00}{#1}}
\newcommand{\RegionMarkerTok}[1]{\textcolor[rgb]{0.00,0.23,0.31}{#1}}
\newcommand{\SpecialCharTok}[1]{\textcolor[rgb]{0.37,0.37,0.37}{#1}}
\newcommand{\SpecialStringTok}[1]{\textcolor[rgb]{0.13,0.47,0.30}{#1}}
\newcommand{\StringTok}[1]{\textcolor[rgb]{0.13,0.47,0.30}{#1}}
\newcommand{\VariableTok}[1]{\textcolor[rgb]{0.07,0.07,0.07}{#1}}
\newcommand{\VerbatimStringTok}[1]{\textcolor[rgb]{0.13,0.47,0.30}{#1}}
\newcommand{\WarningTok}[1]{\textcolor[rgb]{0.37,0.37,0.37}{\textit{#1}}}

\usepackage{longtable,booktabs,array}
\usepackage{calc} % for calculating minipage widths
% Correct order of tables after \paragraph or \subparagraph
\usepackage{etoolbox}
\makeatletter
\patchcmd\longtable{\par}{\if@noskipsec\mbox{}\fi\par}{}{}
\makeatother
% Allow footnotes in longtable head/foot
\IfFileExists{footnotehyper.sty}{\usepackage{footnotehyper}}{\usepackage{footnote}}
\makesavenoteenv{longtable}
\usepackage{graphicx}
\makeatletter
\newsavebox\pandoc@box
\newcommand*\pandocbounded[1]{% scales image to fit in text height/width
  \sbox\pandoc@box{#1}%
  \Gscale@div\@tempa{\textheight}{\dimexpr\ht\pandoc@box+\dp\pandoc@box\relax}%
  \Gscale@div\@tempb{\linewidth}{\wd\pandoc@box}%
  \ifdim\@tempb\p@<\@tempa\p@\let\@tempa\@tempb\fi% select the smaller of both
  \ifdim\@tempa\p@<\p@\scalebox{\@tempa}{\usebox\pandoc@box}%
  \else\usebox{\pandoc@box}%
  \fi%
}
% Set default figure placement to htbp
\def\fps@figure{htbp}
\makeatother





\setlength{\emergencystretch}{3em} % prevent overfull lines

\providecommand{\tightlist}{%
  \setlength{\itemsep}{0pt}\setlength{\parskip}{0pt}}



 


\KOMAoption{captions}{tableheading}
\makeatletter
\@ifpackageloaded{caption}{}{\usepackage{caption}}
\AtBeginDocument{%
\ifdefined\contentsname
  \renewcommand*\contentsname{Table of contents}
\else
  \newcommand\contentsname{Table of contents}
\fi
\ifdefined\listfigurename
  \renewcommand*\listfigurename{List of Figures}
\else
  \newcommand\listfigurename{List of Figures}
\fi
\ifdefined\listtablename
  \renewcommand*\listtablename{List of Tables}
\else
  \newcommand\listtablename{List of Tables}
\fi
\ifdefined\figurename
  \renewcommand*\figurename{Figure}
\else
  \newcommand\figurename{Figure}
\fi
\ifdefined\tablename
  \renewcommand*\tablename{Table}
\else
  \newcommand\tablename{Table}
\fi
}
\@ifpackageloaded{float}{}{\usepackage{float}}
\floatstyle{ruled}
\@ifundefined{c@chapter}{\newfloat{codelisting}{h}{lop}}{\newfloat{codelisting}{h}{lop}[chapter]}
\floatname{codelisting}{Listing}
\newcommand*\listoflistings{\listof{codelisting}{List of Listings}}
\makeatother
\makeatletter
\makeatother
\makeatletter
\@ifpackageloaded{caption}{}{\usepackage{caption}}
\@ifpackageloaded{subcaption}{}{\usepackage{subcaption}}
\makeatother
\usepackage{bookmark}
\IfFileExists{xurl.sty}{\usepackage{xurl}}{} % add URL line breaks if available
\urlstyle{same}
\hypersetup{
  pdftitle={Problem Set 2},
  colorlinks=true,
  linkcolor={blue},
  filecolor={Maroon},
  citecolor={Blue},
  urlcolor={Blue},
  pdfcreator={LaTeX via pandoc}}


\title{Problem Set 2}
\author{}
\date{}
\begin{document}
\maketitle


\begin{enumerate}
\def\labelenumi{\arabic{enumi}.}
\item
  \textbf{Cancer and Group Therapy}. Researchers randomly assigned
  metastatic breast cancer patients to either a control group or a group
  that received weekly 90-minute sessions of group therapy and
  self-hypnosis. The group therapy involved discussion and support for
  coping with the disease. The goal of the experiment was to see whether
  the latter treatment improved the patients' quality of life at the
  time, but a followup study on these patients collected data on the
  number of months of survival after the beginning of the
  study\footnote{Data from Q 4.31 in \emph{Statistical Sleuth}.}.

  You can access the data from this study using the following code.

\begin{Shaded}
\begin{Highlighting}[]
\FunctionTok{library}\NormalTok{(tidyverse)}
\NormalTok{cancer }\OtherTok{\textless{}{-}} \FunctionTok{read.csv}\NormalTok{(}\StringTok{"https://stat158.berkeley.edu/spring{-}2026/data/breast{-}cancer/breast{-}cancer.csv"}\NormalTok{)}
\end{Highlighting}
\end{Shaded}

  \texttt{GROUP} is the original group assigned to the subject.
  \texttt{SURVIVAL} is the survival time in months from the beginning of
  the study. At the time the survival data was collected (10 years
  later), some subjects were still alive. They are flagged in the
  \texttt{CENSOR} column. Data (in this case survival time) that can
  only be known up to some bound is called \emph{censored}.

  \begin{enumerate}
  \def\labelenumii{\alph{enumii}.}
  \tightlist
  \item
    Choose an effective way to visualize the data and present the plot.
    Comment on what you see.
  \item
    Is there evidence of an effect of the group therapy on survival
    time?

    \begin{enumerate}
    \def\labelenumiii{\roman{enumiii}.}
    \tightlist
    \item
      Answer this using randomization inference and two different test
      statistics: the difference in means and the difference in
      medians\footnote{You can augment the \texttt{rand\_stats()}
        function to take
        \texttt{stat\ =\ c("diff\ in\ mean",\ "diff\ in\ medians")} as
        an argumnent and then copy and modify the \texttt{if} block to
        calculate the appropriate statistic when
        \texttt{stat\ ==\ "diff\ in\ medians"}.}.
    \item
      Answer this using model-based inference using a two-sample t-test
      using the pooled estimate of the standard deviation\footnote{This
        should be a review from your course in Statistical Inference. If
        you need a brush up, see
        \url{https://www.itl.nist.gov/div898/handbook/eda/section3/eda353.htm}
        and \texttt{?t.test} in R.}.
    \item
      Compare the methods by first describing the different ways they
      conceive of randomness and then comparing the results and what
      they say about the research question.
    \end{enumerate}
  \item
    When conducting Randomization Based Inference, as you did above,
    what group of units does the inference directly apply to? Said
    another way, the parameter under study is a function of which group
    of units?
  \item
    What problems do you see with having censored data and what problem
    could it \emph{potentially} create for your analysis? Do you think
    it has affected your conclusions for this particular dataset?
  \end{enumerate}
\end{enumerate}

\newpage{}

\begin{enumerate}
\def\labelenumi{\arabic{enumi}.}
\setcounter{enumi}{1}
\item
  \textbf{Statistical Errors in Testing}. Consider the study of whether
  group therapy can have an effect on breast cancer survival.

  \begin{enumerate}
  \def\labelenumii{\alph{enumii}.}
  \tightlist
  \item
    What would constitute a type I error?
  \item
    What would constitute a type II error?
  \item
    Why do statistical errors in testing take place? Is there any way to
    completely eliminate them?
  \end{enumerate}
\end{enumerate}

\begin{enumerate}
\def\labelenumi{\arabic{enumi}.}
\setcounter{enumi}{2}
\item
  \textbf{The Levers of Power}. Consider the study of whether group
  therapy can have an effect on breast cancer survival.

  \begin{enumerate}
  \def\labelenumii{\alph{enumii}.}
  \tightlist
  \item
    What is statistical power in this context?
  \item
    Imagine you were advising the researchers as they were planning
    their study. What are four different parameters that effect the
    power that they could consider optimizing? Realizing that they have
    practical constraints on what they can do, what guidance can you
    give for selecting values of those parameters?
  \end{enumerate}
\end{enumerate}

\newpage{}

\begin{enumerate}
\def\labelenumi{\arabic{enumi}.}
\setcounter{enumi}{3}
\item
  \textbf{Anchoring Power Curve (function of effect size)}. A power
  curve shows, for a particular experiment, the relationship between
  some parameter in the experiment and the statistical power. In
  Randomization-Based Inference, they are approximated point-by-point,
  calculating the power for different values of one parameter while
  keeping the others fixed.

  Create a power curve that shows the relationship between the effect
  size (as measured by \(\tau\), a ``constant shift'' ITE) and the power
  for the Anchoring Experiment. Recall this is a two-tailed test
  conducted at \(\alpha = .05\). Fix the sample size at the values
  observed in the data. Consider at least five values for \(\tau\)
  between 0 and 20 (you're encouraged to consider more values and
  consider negative values). Put those values of \(\tau\) along the
  x-axis and the corresponding values of power on the y-axis.

  \emph{Tip: this exercise will require a lot of code unless you wrap
  stretches of it in functions. See the last slides from class for
  functions that you're welcome to use (if you'd like more practice, try
  writing your own first)}.
\end{enumerate}

\newpage{}

\begin{enumerate}
\def\labelenumi{\arabic{enumi}.}
\setcounter{enumi}{4}
\item
  \textbf{Practice with the t}.\footnote{In R, for the \(t\)
    distribution, the PDF is \texttt{dt()}, the quantile function is
    \texttt{qt()}, and the CDF is \texttt{pt()}. There are analogously
    named functions for all common named probability distributions,
    e.g.~\texttt{dbinom()} and \texttt{pnorm()}.} Create a visualization
  of each of the following functions. Ensure the axes are labelled in an
  informative manner and add appropriate titles. Choose x-axis limits
  that focus on the most interesting part of the functions.\footnote{An
    easy way to plot functions in R is to create a data frame with a
    column \texttt{x} that has a sequence of many evenly-spaced values
    along the x-axis (try 100 values). You can then add additional
    columns for each function of \texttt{x} that you wish to visualize.
    You can plot a pair of those columns against one another and connect
    the points with a line.}

  \begin{enumerate}
  \def\labelenumii{\alph{enumii}.}
  \tightlist
  \item
    The probability density function (PDF) of the \(t\) distribution
    with 4 degrees of freedom with vertical lines marking the .025 and
    .975 quantiles of the distribution.
  \item
    The cumulative distribution function (CDF) of the \(t\) distribution
    with 4 dsegrees of freedom.
  \item
    The PDF of the non-central \(t\) distribution with 4 degrees of
    freedom and a non-centrality parameter of .3.
  \end{enumerate}
\end{enumerate}

\begin{enumerate}
\def\labelenumi{\arabic{enumi}.}
\setcounter{enumi}{5}
\item
  \textbf{Anchoring Model-Based Power Curve (function of effect size)}.

  \begin{enumerate}
  \def\labelenumii{\alph{enumii}.}
  \item
    Using the \(t\) distribution, create a power curve that shows, for
    the Anchoring Experiment, the relationship between the effect size
    (as measured by \(\mu_1 - \mu_0\)) and the power for a two-tailed
    test conducted at \(\alpha = .05\). Fix the sample size at the
    values observed in the data and use 400 as your value of
    \(\sigma^2\).
  \item
    At what effect size does the power reach .8?
  \item
    \emph{Optional Challenge:} Overlay the power curve you created in
    part a with the power curve you via randomization-based inference in
    the previous exercise. Do they look similar? If not, can you explain
    why not?
  \end{enumerate}
\end{enumerate}




\end{document}
